\documentclass[Total.tex]{subfiles}

\begin{document}
	Some of contents:
	\begin{itemize}
		\item Smooth Manifold;
		\item De Rham Cohomology;
		\item Integratiability, Foliation;
		\item Differential Geometry without Metric(Connection Geometry);
		\item 
	\end{itemize}
	
	And the following contents are not contained:
	\begin{itemize}
		\item Specified Geometry;
		\item Characteristic Classes;
		\item Algebraic Topology: Manifolds;
		\item Chern-Weil;
		\item Analysis on Manifolds;
	\end{itemize}
	\chapter{Smooth Manifold}
		\section{Definitions}
			\begin{Def}
				An $n$-dimensional \textbf{smooth manifold} is a topological manifold $(M,\{(U_i,\varphi_i)\})$ with \textbf{smooth transition function}
				\begin{gather*}
					\varphi_i\circ\varphi_j^{-1}\in C^\infty_{diff}(\hat{U}_j,\hat{U}_i)
				\end{gather*}
			\end{Def}
			And for every smooth charts, there exists a \textbf{maximal chart} which contains the smooth charts. Such atlas is called the \textbf{smooth structure} of $M$.
			
			\begin{lemma} Let $M$ be a topological manifold:
				\begin{itemize}
					\item Every smooth atlas is contained in a maximal smooth atlas;
					\item Two smooth atlases for $M$ determines the same maximal smooth atlas iff their union is smooth.
				\end{itemize}
			\end{lemma}
			\begin{proof}
				\begin{itemize}
					\item Let $\mathcal{A}$ be the collection of charts that are smooth compatible with every chart in $\bar{\mathcal{A}}$. 
					
					 $\bar{\mathcal{A}}$ is a smooth chart. Firstly, $\mathcal{A}\subseteq \bar{\mathcal{A}}$. For two charts $(U,\varphi),(V,\psi)\in \bar{\mathcal{A}}$ where $U\cap V\not=\varnothing$, select some $p\in U\cap V$ and a chart $(W,\theta)$ where $p\in U\cap V\cap W$. Omitted.
					\item Omitted.
				\end{itemize}
			\end{proof}
			
			\begin{theorem}
				(Milnor,1956) $S^7$ has more than one the smooth structures.
			\end{theorem}
			\begin{theorem}
				(Exotic $\RR^4$) $\RR^4$ has infinite many of smooth structures.
			\end{theorem}
			
			\subsection{Smooth Functions}
			\begin{Def}
				A function $f:M\rightarrow \RR$ is \textbf{smooth function} if for every coordinate $(U_i,\varphi_i)$
				\begin{gather*}
					f\circ \varphi_i^{-1}:\hat{U}_i\rightarrow \RR
				\end{gather*} 
				is smooth.
			\end{Def}
			\subsection{Construction of Smooth Manifold}
			\begin{lemma}
				(Smooth Manifold Construction Lemma) Let $M$ be a set, and suppose we are given a collection $\{U_\alpha\}$ of subsete of $M$, together with injective maps $\varphi_\alpha:U_\alpha\rightarrow \RR^n$. With following properties:
				\begin{itemize}
					\item For each $\alpha,\varphi_\alpha(U_\alpha)$ is an open subset of $\RR^n$;
					\item $\varphi_\alpha(U_\alpha\cap U_\beta)$ are open in $\RR^n$;
					\item Whenever $U_\alpha\cap U_\beta\not=\varnothing,\varphi_\alpha\circ \varphi_\beta^{-1}$ is a diffeomorphism;
					\item Countably many $U_\alpha$ cover $M$;
					\item Hausdorfness;
				\end{itemize}
			\end{lemma}
			\begin{proof}
				Omitted.
			\end{proof}
		\section{Smooth Functions and Smooth Maps}
			\subsection{Definitions}
			\begin{Def}
				\begin{itemize}
					\item A function $f:M\rightarrow \RR^k$ is \textbf{smooth} if for every $p\in M$, there exists a smooth chart $(U,\varphi)$ and the composition $f\circ \varphi^{-1}:\hat{U}\rightarrow \RR^k$ is smooth;
					
					Such $f\circ \varphi^{-1}$ is called the \textbf{coordinate representation} of $f$;
					\item Furthermore, $F:M\rightarrow N$ is smooth if for some charts $(U,\varphi),(V,\psi)$,
					\begin{gather*}
						\varphi\circ F\circ \psi^{-1}:\psi(U\cap V)\rightarrow \varphi(U\cap V)
					\end{gather*}
					is smooth
				\end{itemize}
			\end{Def}
			
			It is easy to verify:
			\begin{itemize}
				\item $C^\infty(M)$ is commutative $\RR$-algebra;
				\item For every charts, the representation map is smooth;
			\end{itemize}
			
			\begin{lemma}
				(Smoothness is Local)Let $M,N$ be two smooth manifolds, $F:M\rightarrow N$ be a map. If every point $p\in M$, restriction $F|U$ is smooth, then $F$ is smooth; Conversely also holds.
			\end{lemma}
			\begin{proof}
				Omitted;
			\end{proof}
			
			\begin{lemma}
				(Pasting) If given $F_\alpha:U_\alpha\rightarrow N$ with $F_\alpha|U_\alpha\cap U_\beta=F_\beta|U_\alpha\cap U_\beta,\cup U_\alpha=M$. Then there exists $F:M\rightarrow N$ such that
				\begin{gather*}
					F|U_\alpha=F_\alpha
				\end{gather*}
			\end{lemma}
			\begin{proof}
				内容...
			\end{proof}
			
			\begin{lemma}
				Every smooth map between smooth manifolds is continuous.
			\end{lemma}
			\begin{proof}
				Omitted.
			\end{proof}
			
			\subsection{Diffeomorphisms}
			\begin{Def}
				内容...
			\end{Def}
		\section{Smooth Covering Maps}
		\section{Proper Maps}
			\begin{Def}
				$F:M\rightarrow N$ is \textbf{proper} if the inverse image of any compact $K\subseteq N$ is compact.
			\end{Def}
			\begin{lemma}
				Suppose $M$ is a compact space and $N$ is a Hausdorff space, then every continuous map $F:M\rightarrow N$ is proper. 
			\end{lemma}
			\begin{proof}
				内容...
			\end{proof}
		\section{Partition Unity}
	\chapter{Tangent Spaces}
		\section{Tangent Vectors}	
			\subsection{Tangent Space as Derivation of the Space of Germs}
				\subsubsection{Germs}
				\begin{Def}
					\begin{itemize}
						\item 
						The \textbf{germ} on $C_p^\infty(M)$ is specified by the equivalence class by:
						\begin{gather*}
							(U,f)\simeq (V,g)\Leftrightarrow \exists W,p\in W\subseteq U\cap V, f|W=g|W
						\end{gather*}
						
						And operations
						\begin{gather*}
							内容...
						\end{gather*}
						Thus is a linear space;
						\item A \textbf{derivation} of $C_p^\infty(M)$ is a linear map
						\begin{gather*}
							v:C^\infty_p(M)\rightarrow \RR
						\end{gather*}
						
						which satisfies Lebnitz's law:
						\begin{gather*}
							内容...
						\end{gather*}
					\end{itemize}
				\end{Def}
				
			\subsection{Tangent Vectors as Equivalence Classes of Vectors}
				\begin{Def}
					内容...
				\end{Def}
		\section{Differentials}
			\subsection{Derivative on Opens}
				
				\begin{Def}
					内容...
				\end{Def}
				
			\subsection{Formel Derivations}
				\begin{Def}
					Let $M$be a manifold (with or without boundary), then a linear map
					\begin{gather*}
						C^\infty(M)^*\ni v:C^\infty(M)\rightarrow \RR
					\end{gather*}
					
					such that:
					\begin{itemize}
						\item (Lebnitz's Law) $v(fg)=f(p)v(g)+g(p)v(f)$;
					\end{itemize}
				\end{Def}
				
				\begin{proposition}
					The follwoings vector spaces are isomorphic:
					\begin{itemize}
						\item 
						\item 
						\item 
					\end{itemize}
				\end{proposition}
				\begin{proof}
					内容...
				\end{proof}
			\subsection{Differential of Smooth Map}
				\begin{Def}
					内容...
				\end{Def}
				
			\subsection{Pushforwards}
			
		\section{Coordinate Computations}
		\section{Tangent Vectors to Curves,Fibres}
			\subsection{Curves}
			\subsection{Fibres}
		\section{Tangent Bundles}
			\begin{Def}
				Define: For manifold $M$, the \textbf{tangent space}
				\begin{gather*}
					T M=\coprod_{p\in M} T_p M
				\end{gather*}
			\end{Def}
			Then:
			\begin{proposition}
				(Smooth Structure of $TM$, Smoothness of Bundle Map)
			\end{proposition}
			\begin{proof}
				内容...
			\end{proof}
			\subsection{Functoriality}
			
			\begin{proposition}
				$T:Man^1\rightarrow VBun$ is a functor. 
			\end{proposition}
			\begin{proof}
				\begin{itemize}
					\item Firstsly, $T$ is well-defined;
					\item Secondly, $T(g\circ f)=T(g)\circ T(f)$. $f:M\rightarrow N,g:N\rightarrow L$,
					
					\begin{gather*}
						T_p(g\circ f)(X)(h)=X(h\circ (g\circ f))\\
						\qquad =T_p(f)(X)(h\circ g)\\
						\qquad =T_p(g)\circ T_p(f)(X)
					\end{gather*}
					
					where $X\in T_p M$. 
				\end{itemize}
			\end{proof}
			
			Especially
			
			\begin{proposition}
				If $f:M\rightarrow N$ is smooth map, so is $T(f)$. 
			\end{proposition}
			\begin{proof}
				Consider followin diagram
				\begin{center}
					\begin{tikzcd}
						TM\arrow[r,"T(f)"]\arrow[d]&TN\arrow[d]\\
						M\arrow[r,"f"]&N
					\end{tikzcd}
				\end{center}
				
				Take charts 
			\end{proof}
			
			\subsection{Pfaffian Form}
			
			\begin{Def}
				Let $M$ be a differentiable manifold and $V$ a finite dimensional vector space. A \textbf{$V$-valued Pfaffian form}, also called $1$-form on $M$ is a smooth map
				\begin{gather*}
					\omega:TM\rightarrow V
				\end{gather*}
				where $\omega_p:T_p M\rightarrow V$ is linear. 
			\end{Def}
			
			The set of all Pfaffian forms on $M$
			\begin{gather*}
				\Omega^1(M,V)
			\end{gather*}
			
			Especially, when $V=\RR$, denoted by $\Omega_1(M)$. As a matter of fact
			\begin{gather*}
				\Omega^1(M,V)=\{\omega|\omega_p:T_p M\rightarrow V,p\in M\}
			\end{gather*}
			
			\begin{Def}
				Define the product
				\begin{gather*}
					C^\infty(M)\times \Omega^1(M,V)\rightarrow \Omega^1(M,V)\\
					f\cdot \omega:=(f\circ \pi_{TM})\cdot \omega
				\end{gather*}
			\end{Def}
			
			\begin{example}
				内容...
			\end{example}
			
			\subsubsection{Pullback}
			
			\begin{Def}
				Let $f:N\rightarrow M$. Then, the \textbf{pullback}
				\begin{gather*}
					f^*:\Omega^1(M,V)\rightarrow \Omega^1(N,V)\\
					\omega\mapsto \omega\circ (Tf)
				\end{gather*}
			\end{Def}
			
			\subsubsection{Pfaffian Form with Vector Bundles}
				\qquad Consider the vector bundle:
				\begin{gather*}
					\pi:E\rightarrow M\qquad \tau:T^*M\rightarrow M
				\end{gather*}
				
				Notice that:
				\begin{gather*}
					\mbox{(Locally Trivial)}\qquad \pi^{-1}(U)\cong U\times \RR^n
				\end{gather*}
				
				Then, define:
				\begin{gather*}
					\Omega^1(M;E)=\Gamma(T^*M\otimes E)
				\end{gather*}
				
				And we observe the section locally, with local triviality $\pi^{-1}(U)\cong U\times \RR^n$, we have
				\begin{gather*}
					\Gamma(T^*\otimes E)|U=\Omega(U;\RR^n)
				\end{gather*}
			\subsection{Curve, and its Derivative}
			\begin{Def}
				\begin{itemize}
					\item 
					A \textbf{smooth curve} is a smooth function
					\begin{gather*}
						\gamma:I\rightarrow M
					\end{gather*}
					\item  If $\gamma$ could be seen as a smooth map between manifolds $I$ and $M$, now under the standard basis of $I$, we obtain
					\begin{gather*}
						(d\gamma)_t: T_tI=\RR\rightarrow T_{\gamma(t)} M
					\end{gather*}
					
					Especially, with the basis $(\frac{\partial }{\partial t}|t_0)$, define
					\begin{gather*}
						\dot{\gamma}: I\rightarrow TM\qquad t\mapsto (d\gamma)_{t_0}(\frac{\partial}{\partial t}|t_0)\in T_{\gamma(t_0)}M
					\end{gather*}
				\end{itemize}
			\end{Def}
				\begin{proposition}
					\begin{itemize}
						\item 
						We have:
						\begin{gather*}
							\tau \circ \dot{\gamma}=\gamma
						\end{gather*}
						\item And $\dot{\gamma}$ is actually the section of pullback bundle: The pullback bundle induced by $\gamma$:
						\begin{center}
							\begin{tikzcd}
								\gamma^* TM\arrow[r,dashrightarrow]\arrow[d,dashrightarrow]&TM\arrow[d]\\
								I\arrow[r]&M
							\end{tikzcd}
						\end{center}
						
						Then, $\dot{\gamma}\in \gamma^*TM$. 
					\end{itemize}
				\end{proposition}
				\begin{proof}
					\begin{itemize}
						\item Definition;
						\item Omitted;
					\end{itemize}
				\end{proof}
				
		\section{Vector Fields and Tensor Fields}	
			\begin{Def}
				Let $M$ be a manifold. 
				\begin{itemize}
					\item 
					A \textbf{vector field} on $M$ a section of the map $\pi:TM\rightarrow M$.
					
					\item  More concretely, is a continuous map $X:M\rightarrow TM$, with
					\begin{gather*}
						\pi\circ X=Id_M
					\end{gather*}
					
					Denoted by: 
					\begin{itemize}
						\item $\mathfrak{X}(M)$; or
						\item $\Gamma(M,TM)$; or
						\item $H^0(M,\mathcal{T}_M)=\Gamma(M,\mathcal{T}_M)$; For sheaf $\mathcal{T}_M(U):=\Gamma(U,TU)$;
					\end{itemize}
				\end{itemize}
				
				Equivalently, this is a continuous map with $X_p\in T_p M$;
			\end{Def}
			
			\begin{proposition}
				内容...
			\end{proposition}
			\begin{proof}
				内容...
			\end{proof}
			
			\subsubsection{Smoothness of Vector Field}
			
			\begin{proposition}
				(Smooth Vector Fields as $C^\infty(M)$-Module)
			\end{proposition}
			\begin{proof}
				内容...
			\end{proof}
			\subsubsection{Extension of Vector Field}
			
			\subsubsection{Local and Global Frames}
		\subsection{Vector Fields as Derivation of $C^\infty(M)$}
			\begin{Def}
				Define the action of $X\in \mathfrak{X}(M)$:
				\begin{gather*}
					X:C^\infty(M)\rightarrow C^\infty(M)\qquad f\mapsto (Xf:p\mapsto (Xf)(p)=X_p f)
				\end{gather*}
			\end{Def}
			
			\begin{proposition}
				For $X\in \mathfrak{X}(M);f,g\in C^\infty(M)$ we have
				\begin{gather*}
					X(fg)=X(f)g+X(g)f\qquad X_p(fg)=X_p(g)f(p)+X_p(f)g(p)
				\end{gather*}
			\end{proposition}
			\begin{proof}
				内容...
			\end{proof}
			
			\begin{proposition}
				Conversely, given $D\in Der_\RR(C^\infty(M),C^\infty(M))$, then
				\begin{gather*}
					X:M\rightarrow TM\qquad X_p(f):=(Df)(p)
				\end{gather*}
			\end{proposition}
			\begin{proof}
				We should verify:  $X_p$ is the element of $T_p M$. 
			\end{proof}
			
			\begin{corollary}
				There exists an isomrophsim
				\begin{gather*}
					\mathfrak{X}(M)\cong Der_\RR(C^\infty(M),C^\infty(M))
				\end{gather*}
			\end{corollary}
			\begin{proof}
				Omitted.
			\end{proof}
		\subsection{Lie Derivatives}
			\begin{lemma}
				If $U\subseteq M$ is an open subset and $f\in C^\infty(U,V)$ is a smooth function on $U$ with values in $V$, where $V$ is a finite dimensional vector space, $X\in \mathcal{V}(M)$.
				
				Then we obtain for each a smooth function
				\begin{gather*}
					\mathcal{L}_X f:=df\circ X|U\in C^\infty(U,V)
				\end{gather*} 
			\end{lemma}
			\begin{proof}
				Omitted.
			\end{proof}
			
			\begin{Def}
				Such $\mathcal{L}_X(f)$ is called \textbf{Lie derivative} of $f$ with respect to $X$.
			\end{Def}
			\subsection{Lie Bracket}
			\begin{Def}
				Define the \textbf{Lie bracket} on $\mathfrak{X}(M)$: For $X,Y\in \mathfrak{X}(M)$, 
				\begin{gather*}
					[X,Y]:C^1(M)\rightarrow C^\infty(M)\qquad [X,Y]f:=X(Yf)-Y(Xf)
				\end{gather*}
			\end{Def}
			\begin{proposition}
				$(\mathfrak{X}(M),[-,-])$ is a Lie algebra. And $[X,Y]$ is smooth.
			\end{proposition}
			\begin{proof}
				内容...
			\end{proof}
			
			\begin{proposition}
				(Naturality)
			\end{proposition}
			\begin{proof}
				内容...
			\end{proof}
			\subsubsection{Coordinate Calculating of Lie Bracket}
			
		\section{Tangent Covectors on Manifolds}
			\begin{Def}
				Let $M$ be a manifold. For each $p\in M$, we define the \textbf{cotangent space} at $p$, 
				\begin{gather*}
					T_p^*M:=(T_p M)^\vee
				\end{gather*}
			\end{Def}
			\subsection{Conservative Covector Fields}
	\chapter{Bundles}
		Some of contents should refer to the note: Bundles. We will discuss the bundles in general there. Others: Principle Bundles, Associated Bundles, usw. 
		\section{Vector Bundles}
			
			\begin{Def}
				内容...
			\end{Def}
			\subsubsection{$C^\infty(U)$-Module Structure}
				\begin{Def}
					The vector bundle $s\in \Gamma(U)$ has a $C^\infty(U)$-structure: For $f\in C^\infty(U)$,
					\begin{gather*}
						f\cdot s\in C^\infty(U)
					\end{gather*}
				\end{Def}
			\subsection{Vector-Bundle Valued Differential Forms}
				\begin{Def}
					Let $\pi:E\rightarrow M$ be a vector bundle. Now, there exists two bundles
					\begin{gather*}
						\pi:E\rightarrow M\qquad \tau:TM\rightarrow M\Rightarrow \Lambda^k \tau\otimes \pi
					\end{gather*}
					Define:
					\begin{gather*}
						\Omega^k(M;E):=\Gamma(\Lambda^k T^*M\otimes E)
					\end{gather*}
				\end{Def}
				
				Especially: We will not define the differential here, to make $(\Omega^i(M;E))$ a chain complex. 
				
				Because: There exists no differential $d$ without a structure of connection. We will discuss thsi later. 
				\subsubsection{$\Omega^0(M;E)$}
				
				We have $\Omega^0(M,E)=\Gamma(E)$;
				
				\subsubsection{$\Omega^1(M;E)$, $\Omega^k(M;E)$}
				
				\begin{proposition}
					For vector bundle $\pi:E\rightarrow M$:
					\begin{gather*}
						\Omega^1(M;E)=\Gamma(T^*M\otimes E)\cong Hom_{C^\infty(M)}(\mathfrak{X}(M),\Gamma(E))\\
						\Omega^k(M;E)=\Gamma(\Lambda^k T^* M\otimes E)\cong Alt^k_{C^\infty(M)}(\mathfrak{X}(M),\Gamma(E))
 					\end{gather*}
				\end{proposition}
				\begin{proof}
					We just need to verify, that
					\begin{gather*}
						\Gamma(\Lambda^k T^*M\otimes E)\cong Alt^k_{C^\infty(M)}(\mathfrak{X}(M),\Gamma(E))
					\end{gather*}
					
					Because in linear algebra:
					\begin{gather*}
						\Lambda^k V^*\otimes W\cong Alt^k(V,M)
					\end{gather*}
				\end{proof}
				
			\subsubsection{Local Frame Calculating}
			
			
			\subsection{Fibre, and Operations by Fibrewise Construction}
			
			\subsection{Bundle Metric}
				\begin{Def}
					Let $\pi:E\rightarrow M$ be a vector bundle, then, the \textbf{bundle metric} is 
					\begin{gather*}
						k:E\times_M E\rightarrow M\times \RR
					\end{gather*}
					
					where the restriction of $k$ on the fibre of $M$ is a bilinear map over vector space. This implies: 
					\begin{gather*}
						k|\pi^{-1}(m)\times \pi^{-1}(m):k^{-1}(\{m\}\times \RR)\rightarrow \{m\}\times \RR
					\end{gather*}
					
					Because: $E\times_M E=\sqcup_{m\in M} \pi^{-1}(m)\times \pi^{-1}(m),M\times \RR=\sqcup_{m\in M} \RR$
				\end{Def}
				
				Another viewpoint is: That $k:E\otimes_M E\rightarrow \underline{\RR}$, where $\underline{\RR}$ is the trivial vector bundle $M\times \RR\rightarrow M$. This is a bundle morphism.
			
				\begin{lemma}
					Every vector bundle over a paracompact space admits a (positive-definite) metric. 
				\end{lemma}
				\begin{proof}
					We will prove this in the theory of bundles.  
				\end{proof}
				
				\begin{corollary}
					Every bundle on smooth manifold admits a bundle metric. 
				\end{corollary}
				\begin{proof}
					By paracompactness of smooth manifold. 
				\end{proof}
		\section{Tangent Bundles}
		\section{Cotangent Bundles}
		
		
	
		\section{Frame Bundle}
			\begin{Def}
				Let $M$ be an $n$-dimensional manifold. For every $p\in M$, define:
				\begin{gather*}
					Fr_p(TM):=\{(e_1,\dots,e_n)|Span(e_1,\dots,e_n)=T_p M\}
				\end{gather*}
				
				with $Fr(M):=\sqcup_{p\in M} Fr_p (M)$; Proejction $\pi:Fr(M)\rightarrow M$;
			\end{Def}
			
			\begin{lemma}
				As a matter of fact,
				\begin{gather*}
					Fr_p(TM)\cong Iso(\RR^n,T_p M)
				\end{gather*}
			\end{lemma}
			\begin{proof}
				内容...
			\end{proof}
			
			Notice that: $Fr(M)|U\subseteq [\mathfrak{X}(U)]^n$. It contains a Lie-bracket structure, but not over $\RR$.
			
			If $span_\RR(E_1,\dots,E_n)$ is a finite dimensional Lie algebra, then
			\begin{gather*}
				[E_i,E_j]=c_{ij}^kE_k\qquad c_{ij}^k:U\rightarrow \RR
			\end{gather*}
			
			But $c_{ij}^k$ is not constant. If we consider it as $C^\infty(U)$-module, then
			\begin{gather*}
				[E_i,E_j]=c^k_{ij}E_k
			\end{gather*}
			
			is actually a 'Lie bracket'.
			\subsubsection{Einstein's Summation Convention}
				Sometimes, instead of
				\begin{gather*}
					\sum_i a_i
				\end{gather*}
				
				We just use $a_i$, if there is no confusing. 
			\subsection{Principal $GL(n,\RR)$-Bundle Structure}
				
			\begin{proposition}
				For each fibre $Fr_p M$, it contains a $GL_n(\RR)$-transform
				\begin{gather*}
					(e_1,\dots,e_n)\cdot A=(\sum e_j A_{j1},\dots,\sum_j e_j A_{jn})
				\end{gather*}
				
				And such action is: Free and transitive. Thus $\pi$ is a principal $GL_n(\RR)$-bundle. 
			\end{proposition}
			\begin{proof}
				We just need to verify, $\pi$ is locally trivial: Notice that $\pi^{-1}(U)\cong \sqcup_{q\in U} Fr_q(M)\cong U\times GL(n,\RR)$. The isomorphism is given by: Choose a local coordinate $(U;x^1,\dots,x^n)$, with local frame
				\begin{gather*}
					(\frac{\partial}{\partial x^1},\dots,\frac{\partial}{\partial x^n})
				\end{gather*}
				
				Then, we got a section $s:U\rightarrow \sqcup_{p\in U} Fr_p(M),p\mapsto (\frac{\partial}{\partial x^1},\dots,\frac{\partial}{\partial x^n})$, and any frame on $p$ could be represented as 
				\begin{gather*}
					(e_1,\dots,e_n)=(\frac{\partial}{\partial x^i})_{i=1}^n\cdot A
				\end{gather*}
				
				Thus, $Fr(TM)|U\cong GL_n(\RR)\times U$;
			\end{proof}
			The followings are equivalent:
			\begin{itemize}
				\item Local frame;
				\item Section $s:U\rightarrow \sqcup Fr_p(M)$;
				\item Local trivialization;
			\end{itemize}
			
			\begin{proof}
				As a matter of fact: We could take another viewpoint. Given an original local frame
				\begin{gather*}
					s:U\rightarrow Fr(M)|U
				\end{gather*}
				
				Then, we got a form of $s(p)=s_0(p)g(p),p\in U$; Thus we obtain a smooth map
				\begin{gather*}
					g:U\rightarrow GL_n(\RR)
				\end{gather*}
				
				This obtains a bijection:
				\begin{gather*}
					Fr(M)|U\Leftrightarrow C^\infty(U,GL_n(\RR))
				\end{gather*}
				
				Now, for the local trivialization: We choose a local frame $s\in Fr(M)|U$, then corresponds $C^\infty(U,GL_n(\RR))$.
				
				For selected point $q\in U$, and every $u\in Fr(T_p M)$, there exists decomposition
				\begin{gather*}
					u=s(p)g
				\end{gather*}
				
				Thus, $u\mapsto (p,g)\qquad Fr(T_p M)\cong (p,GL_n(\RR))\Rightarrow Fr(TM)\cong \sqcup_{p\in U} (p,GL_n(\RR))=U\times GL_n(\RR)$
			\end{proof}
			
			So, the isomorphism is not canonically. It depends on choice of local frame. 
			
			\begin{proposition}
				Furthermore, there exists associated bundle structure
				\begin{gather*} 
					TM\cong Fr(M)\times_{GL_n(\RR)} \RR^n
				\end{gather*}
			\end{proposition}
			\begin{proof}
				
			\end{proof}
			
			\begin{example}
				(Frame Bundle of $\RR^n$) Especially, the frame bundle of $\RR^n$ is trivial:
				\begin{gather*}
					Fr(\RR^n)\cong \RR^n\times GL_n(\RR)
				\end{gather*}
				
				Because $\RR^n$ has a standard basis on each tangent space. 
			\end{example}
			
			\begin{example}
				(The Action of $SO(3)$ on $TS^2$)
			\end{example}
			
			For the orthogonal frames, we will discuss the in Riemann geometry;
			
			\subsubsection{Geometric Structure}
			
			\begin{itemize}
				\item Riemannian Metrics$\Leftrightarrow$ $O(n)$;
				\item Oriented Riemann Manifolds$\Leftrightarrow SO(n)$;
				\item $GL(m;\CC)\subseteq GL(2m,\RR)$: Almost Complex Structure;
				\item $U(m)$: Almost Hermitian Geometry;
				\item \dots
			\end{itemize}
			
			We will research $G$-structure in other topics.
			
			\subsubsection{Transition Function}
			\begin{Def}
				Now, consider two elements in frame:$p\in U\cap V$, then for $s\in Fr(M)|U,s'\in Fr(M)|V$, we have
				\begin{gather*}
					s|U\cap V=s'|U\cap V\cdot (g_{ij})
				\end{gather*}
				
				Such $(g_{ij}):U\cap V\rightarrow GL_n(\RR)$ is called \textbf{transition function}
			\end{Def}
			
			\subsubsection{Coframe}
			\begin{Def}
				Define the \textbf{coframe}: Given $(E_1,\dots,E_n)\in Fr(M)|U$, then the \textbf{dual coframe} is the unqiue basis:
				\begin{gather*}
					(\theta^1,\dots,\theta^n)\qquad \theta^i\in \Omega^1(U)
				\end{gather*}
				
				such that $\theta^i(E_j)=\delta_{ij}$. 
				
				Then, the \textbf{coframe} on $U$ is
				\begin{gather*}
					Fr(T^*M)
				\end{gather*}
			\end{Def}
			\subsection{Moving Frame}
				\begin{Def}
					\begin{itemize}
						\item 
						A \textbf{moving frame} is a local section of frame bundle
						\begin{gather*}
							s:U\rightarrow Fr(M)
						\end{gather*}
						\item The \textbf{structure constant} of moving frame is $(c_{ij}^k)$ given by:
						\begin{gather*}
							[E_i,E_j]=\sum c^k_{ij}E_k
						\end{gather*}
						
						where the Lie bracket is well-defined on $\mathfrak{X}(U)$;
					\end{itemize}
				\end{Def}
				\subsubsection{Moving Coframe}
				\begin{proposition}
					(Structure Equation of Coframe)Now, consider moving frame $(E_1,\dots,E_n)$, with structure constant
					\begin{gather*}
						[E_i,E_j]=\sum c_{ij}^kE_k
					\end{gather*} ,then for the dual coframe $(\theta^1,\dots,\theta^n)$, we have a basis $\{\theta^i\wedge \theta^j\}_{i<j}$ for the local frame $\Lambda^2 T_* M$: Thus for every $d\theta^k\in \Omega^2(U)$,
					\begin{gather*}
						d\theta^k=-\frac{1}{2}c_{ij}^k \theta^i\wedge \theta^j
					\end{gather*}
					
					Conversely, we could also obtain the structure constant from .
				\end{proposition}
				\begin{proof}
					内容...
				\end{proof}
				\subsubsection{Moving Frame with Connection}
		\section{Orientation Bundle}
			\begin{Def}
				内容...
			\end{Def}
			\section{Connections without Metric Structures}
		\begin{Def}
			Let: $\pi:E\rightarrow M$ be a smooth vector bundle. 
			\begin{itemize}
				\item 
				A \textbf{(Koszul) connection} is a linear map
				\begin{gather*}
					\triangledown:\Gamma(E)\rightarrow \Omega^1(M;E)=\Gamma(T^*M\otimes E)\cong \sqcup_p Hom_\RR (T_p M,E_p)
				\end{gather*}
				
				which satisfies Lebnitz's Law
				\begin{gather*}
					\triangledown(fs)=df\otimes s+f\triangledown s\qquad f\in C^\infty(M),s\in \Gamma(E)
				\end{gather*}
				
				Thus, the value of $\triangledown s,s\in \Gamma(E)$　could be seen as a morphism between vector bundles
				\begin{gather*}
					\triangledown s:TM\rightarrow E
				\end{gather*}
				\item Or equivalently, in \textbf{covariant derivative}
				\begin{gather*}
					\triangledown_{(-)}(-):\mathfrak{X}(M)\times \Gamma(E)\rightarrow \Gamma(E)
				\end{gather*}
				
				It is actually a bilinear form, which satisfies:
				\begin{itemize}
					\item $\triangledown_{fX+gY}s=f\triangledown_X s+g\triangledown_Y s$;(linear)
					\item $\triangledown_X(s+t)=\triangledown_X(t)+\triangledown_X(t)$;(Additive)
					\item $\triangledown_X(fs)=X(f)s+f\triangledown_X s$;(Lebnitz)
				\end{itemize}
				
				Notice, it is not bilinear, and
				\begin{gather*}
					\triangledown_X(-):\Gamma(E)\rightarrow \Gamma(E)\qquad \triangledown_X\in \operatorname{Diff}^1_\RR(\Gamma(E))\\
					\triangledown_{(-)}(s):\mathfrak{X}(M)\rightarrow \Gamma(E)\qquad s\in \Gamma(E)
				\end{gather*}
			\end{itemize}
			
			The two definition corresponds through:
			\begin{gather*}
				(\triangledown s)(X)=\triangledown_X s\qquad X\in \mathfrak{X}(M)=TM\qquad s\in \Gamma(E)\qquad X\in \mathfrak{X}(M)
			\end{gather*}
			
			\end{Def}
		
		The followings are equivalent:
		\begin{itemize}
			\item $\Omega^1(M;E)$;
			\item $\Gamma(T^*M\otimes E)$;
			\item The morphisms $Mor_{VectBund}(TM,E)$;
		\end{itemize}
		
		
		And: $\triangledown_X(-)$ is actually a differential operator: We will learn this in \underline{analysis on manifold}, and towards \underline{geometric analysis}.
		
		\begin{proposition}
			The map
			\begin{gather*}
				\mathfrak{X}(M)\hookrightarrow Diff_\RR^1(\Gamma(E))\qquad X\mapsto \triangledown_X
			\end{gather*}
			
			is injective. 
		\end{proposition}
		\begin{proof}
			Notice that: For every $f\in C^\infty(M)$, $s\in \Gamma(E)$, 
			\begin{gather*}
				0=(\triangledown_X-\triangledown_Y)(fs)\\
				=(X(f)s+f\triangledown_X s)-(Y(f)s+f\triangledown_Y s)
			\end{gather*}
			
			Especially, when $f=1_M$, $(\triangledown_X-\triangledown_Y)(1_M)=0$, thus
			\begin{gather*}
				=(X(f)-Y(f))s
			\end{gather*}
			
			$s$ is arbitrary, thus $X(f)=Y(f)\Rightarrow X=Y$.
		\end{proof}
		
		\subsubsection{Local Frame Calculating}
			\begin{proposition}
				Let $\pi:E\rightarrow M$ be a vector bundle. If $(e_1,\dots,e_r)\in Fr(M)|U$ is a frame on $U$, with dual frame $(\theta^1,\dots,\theta^r)$, then for connection $\triangledown:\Gamma(E)\rightarrow \Omega^1(M;E)$,
				\begin{gather*}
					\triangledown e_j=\sum_i \omega^i_j \otimes e_i\qquad 
				\end{gather*}
				
				where $\omega^i_j\in \Omega^1(U)=\Gamma(T^*U)$ 1-form, and 
				\begin{gather*}
					\omega^i_j=\sum_k \Gamma^{i}_{jk} \theta^k\\
					\Rightarrow \triangledown e_j=\sum_i \sum_k \Gamma^i_{jk} \theta^k\otimes e_i
				\end{gather*}
			\end{proposition}
			\begin{proof}
				Omitted.
			\end{proof}
			
			\subsubsection{Dual Connection}
				\begin{Def}
					Consider a connection: For bundle $E\rightarrow M$, 
					\begin{gather*}
						\triangledown:\Gamma(E)\rightarrow \Omega^1(M;E)
					\end{gather*}
					
					Then, the induced connection:
					\begin{itemize}
						\item  With pairing 
						\begin{gather*}
							<-,->:E\otimes E^*\rightarrow M\times \RR\\
							(e,e')\mapsto (p,<e,e'>)\qquad e\in E_p,e'\in E_p^*
						\end{gather*}
						
						We got
						\begin{gather*}
							<-,->:\Gamma(E\otimes E^*)=\Gamma(E)\otimes \Gamma(E^*)\rightarrow C^\infty(M)
						\end{gather*}
						
						\item Then, define the \textbf{dual connection}
						\begin{gather*}
							\triangledown:\Gamma(E^*)\rightarrow \Omega^1(M;E^*)=\Gamma(T^*M\otimes E^*)\\
							\alpha\mapsto (s\mapsto X(\alpha(s))-\alpha(\triangledown_X s))
						\end{gather*}
						
						It is the \underline{unique} connection such that
						\begin{gather*}
							X(<\alpha,s>)=<\triangledown_X \alpha s>+<\alpha,\triangledown_X s>
						\end{gather*}
						
						Equivalently,
						\begin{gather*}
							d<\alpha,s>=(\triangledown \alpha,s)=<\alpha,\triangledown s>
						\end{gather*}
					\end{itemize}
				\end{Def}
				
				\begin{proposition}
					For the , the followings are equivalent:
					\begin{itemize}
						\item $X(<\alpha,s>)=<\triangledown_X \alpha s>+<\alpha,\triangledown_X s>$;
						\item $d<\alpha,s>=(\triangledown \alpha,s)=<\alpha,\triangledown s>$;
					\end{itemize}
				\end{proposition}
				\begin{proof}
					内容...
				\end{proof}
				
				\begin{proposition}
					The dual connection satisfies the required condition.
				\end{proposition}
				\begin{proof}
					内容...
				\end{proof}
				
			\subsubsection{Pullback of Connection, Covariant Derivative}
				\begin{Def}
					Given:
					\begin{itemize}
						\item Vector bundle $\pi:E\rightarrow M$;
						\item Associated connection and covariant derivative 
						\begin{gather*}
							\triangledown:\Gamma(E)\rightarrow \Omega^1(M;E)\qquad \triangledown_{(-)}^{(-)}:\mathfrak{X}(M)\times \Gamma(E)\rightarrow \Gamma(E)
						\end{gather*}
						\item Smooth map $f:N\rightarrow M$
					\end{itemize}
					
					Then, the \textbf{pullback} is given by: With induced bundle $f^*\pi:f^*E\rightarrow N$ and induced 
					\begin{gather*}
						f^*\triangledown:\Gamma(f^*E)\rightarrow \Omega^1(N;f^*E)\qquad f^*s\mapsto f^*(\triangledown s)\qquad s\in \Gamma(E)
					\end{gather*}
					
					where $f^*:\Omega^1(M;E)=\Gamma(T^*M\otimes E)\rightarrow \Omega^1(N,f^*E)=\Gamma(T^*N\otimes f^*E),\omega\otimes s\mapsto f^*\omega\otimes f^*s$
				\end{Def}
				
				\begin{proposition}
					(Universal Property of Pullback) The following diagram commutes:
					\begin{center}
						\begin{tikzcd}
							\Gamma(E)\arrow[d,"\triangledown"]\arrow[r,"f^*"]&\Gamma(f^*E)\arrow[d,"f^*\triangledown",dashrightarrow]\\
							\Omega^1(M;E)\arrow[r,"f^*"]&\Omega^1(N,f^*E)
						\end{tikzcd}
					\end{center}
					
					
				\end{proposition}
				\begin{proof}
					Omitted;
				\end{proof}
			\subsubsection{Parallel Transport}
				\begin{Def}
					Let
					\begin{itemize}
						\item $\pi:E\rightarrow M$ be a bundle;
						\item Connection, covariant derivative $\triangledown$;
						\item Path $\gamma:I=[a,b]\rightarrow M$, which is smooth;
					\end{itemize}
					
					Then, we obtain:
					\begin{itemize}
						\item Pullback bundle $\gamma^*\pi:\gamma^*E\rightarrow [a,b]$;
						\item Pullback of connection $\gamma^*\triangledown:\Gamma(\gamma^*E)\rightarrow \Omega^1(I,\gamma^*E)$;
					\end{itemize}
					
					Notice, $\Gamma(\gamma^*E)\ni V:I\rightarrow \gamma^*E$, and $TI$ has a global frame $\frac{\partial}{\partial t}$, $T^*I$ for $\partial t\in T^*I=\Gamma(T^*I\otimes \gamma^*E)$.
					
					Now, define the \textbf{covariant derivative}
					\begin{gather*}
						\frac{D^\triangledown}{dt}:\Gamma(\gamma^*E)\rightarrow \Gamma(\gamma^*E)\qquad V\mapsto (\gamma^*\triangledown V)(\partial t)
					\end{gather*}

					\begin{itemize}
						\item $V$ is called \textbf{parallel along} $\gamma$, if $\frac{D^\triangledown V}{dt}=0$.
					\end{itemize}
				\end{Def}
				
				Namely, we should solve the problem in ODE:
				\begin{gather*}
					\begin{cases}
						\frac{D^\triangledown}{dt}V=0\\
						V(a)=v
					\end{cases}
				\end{gather*}
				
				
				\begin{proposition}
					(Turn into Lienar ODE under Local Frame) 
					\begin{itemize}
						\item Under local frame $(e_1,\dots,e_n)$, if 
						\begin{gather*}
							内容...
						\end{gather*}
						
						Then, the parallel equation become:
						\begin{gather*}
							\dot{v}^i+\omega^i_j(\dot{\gamma})v^j=0
						\end{gather*}
						\item And such $v(t)$ is unique. 
					\end{itemize}
				\end{proposition}
				\begin{proof}
					\begin{itemize}
						\item 
						\item By ODE uniqueness theorem;
					\end{itemize}
				\end{proof}
				
				
				
				\begin{Def}
					Given $v\in E_{\gamma(a)}$, there exists a unique parallel section $V:I\rightarrow \gamma^*E$ such that $V(a)=v$.
					
					A \textbf{parallel transport} is a linear transformation
					\begin{gather*}
						P^\triangledown_\gamma:E_{\gamma(a)}\rightarrow E_{\gamma(b)}\qquad P^\triangledown_\gamma(v)=v
					\end{gather*}
				\end{Def}
				
				With these preparation, we will discuss \underline{holonomy group} later. 
				
				\textbf{Remark} More general theory, will be given in \underline{fibre bundle} and \underline{Ehresmann connection}. 
			\subsubsection{Connection in First Jet Exact Sequences}
				\begin{proposition}
					The followings are equivalent:
					\begin{itemize}
						\item A connection $\triangledown:\Gamma(E)\rightarrow \Omega^1(M;E)$;
						\item The first jet exact sequence splits
						\begin{gather*}
							0\rightarrow T^*M\otimes E\rightarrow J^1E\rightarrow E\rightarrow 0
						\end{gather*}
					\end{itemize}
				\end{proposition}
				\begin{proof}
					内容...
				\end{proof}
		\subsection{Principal Bundle from Connection, Pincipal Connection Form}
			\begin{Def}
				Let $\pi:E\rightarrow M$ be a vector bundle with $rank(\pi)=r$, the Koszul connection $\triangledown:\Gamma(E)\rightarrow \Omega^1(M,E)$
				
				Under frame bundle $P=Fr(E)$
				\begin{gather*}
					P\rightarrow M
				\end{gather*}
			\end{Def}
		\subsection{Linear/Affine Connections}
		
		\begin{Def}
			\begin{itemize}
				\item Especially, if $E=TM$, then the connection
				\begin{gather*}
					\triangledown:\Gamma(TM)\rightarrow \Gamma(T^*M\otimes TM)=\Gamma(End_\RR(TM))\cong End_{C^\infty(M)}(\mathfrak{X}(M))
				\end{gather*}
				
				is called \textbf{linear connection}.
				
				\item Under this definition, the covariant derivative: $\Gamma(TM)=\mathfrak{X}(M)$, 
				\begin{gather*}
					\triangledown_{-}(-):\mathfrak{X}(M)\times \mathfrak{X}(M)\rightarrow \mathfrak{X}(M)
				\end{gather*}
				
			\end{itemize}
		\end{Def}
		
		For each $A\in \Gamma(TM\otimes T^* M)$, we have expansion under the frame and dual coframe
		\begin{gather*}
			(e_1,\dots,e_n)\qquad (\theta^1,\dots,\theta^n)
		\end{gather*}
		
		Then, $A=\sum A^i_j e_i\otimes \theta^j$;
		
		Furthermore, a moving frame turns a tensor field in a scalar coefficiet functions. The example above is $(1,1)$ tensor field.
		
		
		\begin{corollary}
			For fixed bundle $\pi:E\rightarrow M$, the collection of linear collections $Conn(E)$ is an affine space, and a correspondent vector space could be:
			\begin{gather*}
				\Omega^1(M;End(E))
			\end{gather*}
		\end{corollary}
		\begin{proof}
			内容...
		\end{proof}
		
		\textbf{Remark} The differential 
		\begin{gather*}
			内容...
		\end{gather*}
		
		is not a (linear) connection, but actually $C^\infty(M)$-linear:
		\begin{gather*}
			A=
		\end{gather*}
		
		\subsubsection{Torsion Tensor}
		\begin{Def}
			Define 
			\begin{itemize}
				\item the \textbf{torsion}
				\begin{gather*}
					T^\triangledown(X,Y):=\triangledown_X Y-\triangledown_Y X-[X,Y]\qquad \mathfrak{X}(M)\times \mathfrak{X}(M)\rightarrow \mathfrak{X}(M)
				\end{gather*}
				\item the \textbf{curvature}
				\begin{gather*}
					R(X,Y)Z=\triangledown_X\triangledown_Y Z-\triangledown_Y\triangledown_XZ-\triangledown_{[X,Y]}Z\qquad \mathfrak{X}(M)\times \mathfrak{X}(M)\times \mathfrak{X}(M)\rightarrow \mathfrak{X}(M)
				\end{gather*}
			\end{itemize}
		\end{Def}
		
		We will calculate the coordinate form of the two valuables in next secion. 
		
		\begin{proposition}
			(Torsion) For the connection $\triangledown$:
			\begin{itemize}
				\item It is anti-symmetric:
				\begin{gather*}
					内容...
				\end{gather*}
				\item Thus, it can be factorized in 
				\begin{gather*}
					T\in \Gamma(\Lambda^2 T^* M\otimes TM)=\Omega^2(M;TM)
				\end{gather*}
			\end{itemize}
		\end{proposition}
		\begin{proof}
			内容...
		\end{proof}
		
		\begin{proposition}
			(Curvature) For the connection $\triangledown$:
			\begin{itemize}
				\item $R(X,Y)=-R(X,Y)$;
				\item $R\in \Gamma(\Lambda^2 T^*M\otimes T^* M\otimes TM)=\Omega^2(M;End(TM))$;
			\end{itemize}
		\end{proposition}
		\begin{proof}
			内容...
		\end{proof}
		
		\begin{Def}
			\begin{itemize}
				\item If $T(X,Y)=0$, then \textbf{torsion-free};
				\item If $R(X,Y)=0$, then \textbf{connection flat};
			\end{itemize}
		\end{Def}
		\subsubsection{Cartan's Structure Theorem I II and Component Expansion}
			\begin{Def}
				Given a frame bundle $(e_1,\dots,e_n)$ of $TM\rightarrow M$ with coframe $(\theta^1,\dots,\theta^n)$, we have 
				\begin{gather*}
					\triangledown e_j=\omega^i_j\otimes e_i\qquad \omega^i_j\in \Gamma(T^*M)=\Omega^1(U)
				\end{gather*}
				
				Then, $\omega=(\omega^i_j)\in \Omega^1(U;\mathfrak{gl}_n(\RR))$;
			\end{Def}
				\begin{theorem}
					(Cartan) For connection $\triangledown$ and torsion $T^\triangledown$, under the frame $(e_1,\dots,e_n)$, and dual frame $(\theta^1,\dots,\theta^n)$, we have
					\begin{gather*}
						T=\sum_i T^i\otimes e_i\\
						T^i=d\theta^i+\sum_j \omega^i_j\wedge \theta^j
					\end{gather*}
				\end{theorem}
				\begin{proof}
					内容...
				\end{proof}
				\begin{theorem}
					(Cartan) For connection $\triangledown$ and curvature $R$, under the frame $(e_1,\dots,e_n)$ and $(\theta^1,\dots,\theta^n)$,
					\begin{gather*}
						R(X,Y)e_j=\sum_i (\Omega^i_j (X,Y))e_i\qquad \Omega^i_j=\theta^i(R(-,-)e_j)\\
						\Omega^i_j=\sum_k d\omega^i_j+\omega^i_k\wedge \omega^k_j
					\end{gather*}
				\end{theorem}
				\begin{proof}
					内容...
				\end{proof}
		\subsection{Existence of a Connection}	
			Recall:
			\begin{lemma}
				Every paracompact manifold has a unity of partion. 
			\end{lemma}
			\begin{proof}
				In 'general topology'
			\end{proof}
			
			\begin{theorem}
				(Existence of Connection) For smooth bundle $\pi:E\rightarrow M$, and $M$ is a paracompact smooth manifold, then:
				
				Every smooth vector bundle admits a connection;
			\end{theorem}
			\begin{proof}
				内容...
			\end{proof}
			
			\begin{corollary}
				The collection of connections $Conn(E)$ is an affine space. i.e. for $A':=\triangledown'-\triangledown$, it is $C^\infty(M)$-linear. 
			\end{corollary}
			\begin{proof}
				内容...
			\end{proof}
			Especially, with condition for Riemann metric $g$, we have
			\begin{gather*}
				\triangledown g=0\qquad T^\triangledown =0
			\end{gather*}
			
			We got Levi-Civita connection.
		\subsection{Flat Connections, and Twisted De Rham Cohomology}
		\subsubsection{Parallel Frame}
			\begin{Def}
				Let $\pi:E\rightarrow M$ be a vector bundle, with connection $\triangledown$:
				\begin{itemize}
					\item A local frame $(e_1,\dots,e_n)$ is called a \textbf{parallel frame}, if
					\begin{gather*}
						\triangledown e_i=0
					\end{gather*}
				\end{itemize}
			\end{Def}
			
			\begin{proposition}
				内容...
			\end{proposition}
			\begin{proof}
				内容...
			\end{proof}
			
		\subsubsection{Flat Connections}
		\begin{Def}
			Let $E\rightarrow M$ be a vector bundle. $\triangledown$ is a connection.
			
			$\triangledown$ is called a \textbf{flat connection}, if for any $X,Y\in \mathfrak{X}(M)$ 
			\begin{gather*}
				内容...
			\end{gather*}
		\end{Def}
			\subsubsection{Covariant Derivative}
			\begin{Def}
				Define: The \textbf{covariant exterior derivative}
				\begin{gather*}
					d_\triangledown:\Omega^k(M;E)\rightarrow \Omega^{k+1}(M;E)\qquad \alpha\otimes s\mapsto d\alpha \otimes s+(-1)^k \alpha\wedge \triangledown s
				\end{gather*}
			\end{Def}
			
			The following rule
			\begin{gather*}
				内容...
			\end{gather*}
			
			is called \textbf{graded Lebnitz rule}.
			
			\textbf{Remarks} Sometimes operator $d_\triangledown$ could be denoted as $D$, if the context is in Cartan's equation. 
			
			\begin{proposition}
				\begin{itemize}
					\item (Derivative)
					\item (Curvature)
				\end{itemize}
			\end{proposition}
			\begin{proof}
				内容...
			\end{proof}
			
			\begin{proposition}
				(Affine Case of Covariant Exterior Derivative, Torsion)
			\end{proposition}
			\begin{proof}
				内容...
			\end{proof}
			\subsubsection{Twisted De Rham Complex}
			\begin{proposition}
				\begin{itemize}
					\item When $k=0$, $d_\triangledown s=\triangledown s$; Thus, $d_\triangledown$ is a continuation of connection;
					\item If $\triangledown$ is flat, then
					\begin{gather*}
						d_\triangledown^2=0
					\end{gather*}
					
					thus $(\Omega^k(M;E))$ is a cochain complex.
				\end{itemize}
			\end{proposition}
			\begin{proof}
				内容...
			\end{proof}
			
			
			\subsubsection{Twisted De Rham Cohomology}
			\begin{Def}
				
			\end{Def}
			
			\subsection{Holonomy Groups}
				\begin{Def}
					Let $\pi:E\rightarrow M$ be a vector bundle with connection $\triangledown$. And fix a point $p\in M$:
					\begin{itemize}
						\item  Consider the loop space with point $(\Omega M,p)$
						
						Define:
						\begin{gather*}
							Hol_p(\triangledown):=\{P_\gamma|\gamma\in \Omega M,P_\gamma\mbox{ is Parallel}\}\subseteq GL(E_p)
						\end{gather*}
						
						is called the \textbf{holonomy group} at point $p$. 
					\end{itemize}
				\end{Def}
				
				\begin{proposition}
					It is actually a group, with
					\begin{gather*}
						P_{\gamma_1}\circ P_{\gamma_2}=P_{\gamma_1\circ \gamma_2}
					\end{gather*}
				\end{proposition}
				\begin{proof}
					内容...
				\end{proof}
				
				\subsubsection{Restricted Holonomy Group}
				\subsubsection{Infinitesimal Holonomy}
				\subsubsection{Flatness}
					\begin{corollary}
						(Flatness) If $R^\triangledown=0$, then 
						\begin{gather*}
							Hol^0_p(\triangledown)=1
						\end{gather*}
					\end{corollary}
					\begin{proof}
						内容...
					\end{proof}
					
				\subsubsection{Ambrose–Singer theorem}
			\subsection{Mappings of Connections}
				\begin{Def}
					内容...
				\end{Def}
			\subsection{Invariant Connections}
		\section{Curvature}
		
		In the pasting discussion, we have know: The differential map of chain complex relies on the flatness of connections. But if not, we have
		\begin{gather*}
			d_{\triangledown}^2\not=0
		\end{gather*}
		
		We could do some research of the obstruntion to the connection-induced differential squaring to zero:
		\begin{Def}
			Suppose $\pi:E\rightarrow M$ be a vector bundle, then, 
			\begin{gather*}
				F_{\triangledown}:=d_\triangledown^2
			\end{gather*}
			
			is called \textbf{curvature}. 
		\end{Def}
		\subsubsection{Affine Connection Case}
		\begin{proposition}
			Especially: For the bundle $\tau:TM\rightarrow M$, we have
			\begin{gather*}
				R=F_\triangledown
			\end{gather*}
		\end{proposition}
		\begin{proof}
			内容...
		\end{proof}
		
		\subsection{Bianchi Identity I II}
			\subsubsection{Bianchi I}
				\begin{theorem}
					(Bianchi) Suppose $E=TM$, $\triangledown$ is an affine connection. 
				\end{theorem}
				\begin{proof}
					内容...
				\end{proof}
			\subsubsection{Bianchi II}
				\begin{theorem}
					(Bianchi)
				\end{theorem}
				\begin{proof}
					内容...
				\end{proof}
				
		\section{Jet Bundles}
		\textbf{Remark} And in the future, we will work with Chern-Weil theory.
	\chapter{Submersions, Immersions and Embeddings}
		\section{Maps of Constant Rank}
		\section{Embeddings}
			\begin{Def}
				Let $M,N$ be two manifolds. Then, a \textbf{smooth embedding} $i:N\rightarrow M$ satisfies:
				\begin{itemize}
					\item $i$ is topological embedding;
					\item $i$ is smooth;
				\end{itemize}
			\end{Def}
			\subsection{Normal Bundle}
				\begin{Def}
					Suppose $i:N\hookrightarrow M$ be a smooth embedding, then define the \textbf{normal bundle}
					
					\begin{gather*}
						v_{N/M}:=i^*TM/TN
					\end{gather*}
					
					Induced by $i^*:TN\rightarrow i^*TM$. The bundle structure
					\begin{gather*}
						v_{N/M}\rightarrow N
					\end{gather*}
				\end{Def}
				
				\textbf{Remark} The structure of normal bundle here depends on the differential structure(Especially, tangential space), thus we could not generalize this definition to the topological manifold. 
				
				An alternative is \textbf{microbundle} given by Milnor. We will introduce this in topological manifold.
				
				\begin{proposition}
					There exists a short exact sequence
					\begin{gather*}
						0\rightarrow TN\stackrel{i^*}{\rightarrow} i^* TM\stackrel{p}{\rightarrow}v_{N/M}\rightarrow 0
					\end{gather*}
					
					and the sequence splits. i.e. There exists $\iota:v_{N/M}\rightarrow i_*TM$ with $p\circ \iota=id_{v_{N/M}}$. 
				\end{proposition}
				\begin{proof}
					We just observe the sequence fibrewise. 
				\end{proof}
				
				Especially, according to our definition of bundle metric, by choosing a positive-definite metric, we obtain 
				\begin{gather*}
					v_{N/M}\cong TN^\perp
				\end{gather*}
				
				Thus, $i_* TM\cong v_{N/M}\oplus TN$. 
				
				\textbf{Remark} This property depends on $Ext^1$ for the nontrivial extension.(If trivial, then splits, and isomorphic to direct sum) 
				
				\begin{proposition}
					(Embedded in $\RR^n$) Suppose $M\subseteq \RR^n$, then $NM$ is an embedded submanifold of 
					\begin{gather*}
						T\RR^n\cong \RR^n\times \RR^n
					\end{gather*}
				\end{proposition}
				\begin{proof}
					内容...
				\end{proof}
				
				\subsection{Tubular Neighborhood}
				Instead of $v_{M/N}$, we will use: $v_{M/N}$ as the map and $Tot(v_{M/N})=Tot(v_i)$ as the total space, $N$ for Base space. 
				
				\begin{theorem}
					Suppose $i:N\hookrightarrow M$ be a smooth embedding. Then, there exists a natural zero-section 
					\begin{gather*}
						0_{N}:N\rightarrow Tot(v_{N/M})=Tot(v_i)\qquad p\mapsto 0_p
					\end{gather*}
					
					Then, there exists an open neighborhood $V\subseteq Tot(v_i)$ of $0_N$, and smooth embedding $\Phi:V\rightarrow M$ such that
					\begin{gather*}
						\Phi\circ 0_N=i
					\end{gather*}
				\end{theorem}
				
				i.e. There exists factorization
				\begin{center}
					\begin{tikzcd}
						N\arrow[rd,hookrightarrow]\arrow[r,"0_N"]&\exists V\arrow[d,"\Phi",hookrightarrow,dashrightarrow]&\subseteq Tot(v_{N/M})\\
						&M
					\end{tikzcd}
				\end{center}
				
				And $\Phi:V\rightarrow W=\Phi(V)\subseteq M$ is a diffeomorphism. 
				
				\begin{proof}
					内容...
				\end{proof}
				
				\begin{corollary}
					If $M\subseteq N$, then there exists 
				\end{corollary}
				\begin{proof}
					内容...
				\end{proof}
				
				Now, notice that $V$ is open in $Tot(v_{N/M})$. Then, it could be described as
				\begin{gather*}
					V=\{(x,v)||v|<\epsilon(x)\}\subseteq Tot(v_{N/M})=\sqcup_{x\in N} i^* T_{i(x )} M/T_xN
				\end{gather*}
				
				Especially, if $N$ is compact, then we obtain that
				\begin{gather*}
					V=\{(x,v)||v|<\epsilon\}
				\end{gather*}
		\section{Submersions}
		
		\section{Smooth Covering Maps}
			Recall
			\begin{Def}
				内容...
			\end{Def}
	\chapter{Submanifolds}
		\section{Embedded Submanifold}
			\begin{proposition}
				Suppose $F:M\rightarrow N$ smooth, $q\in N$ regular value. Then, for $p\in f^{-1}(q)$,
				\begin{gather*}
					T_p(f^{-1}(q))=Ker(d_p F)
				\end{gather*}
			\end{proposition}
			\begin{proof}
				内容...
			\end{proof}
		\subsection{Immersed Submanifold}
			\begin{Def}
				Immersed submanifold: Image of injective immersion.
			\end{Def}
			\begin{example}
				$\RR\rightarrow T,x\mapsto (exp(2\pi ix),exp(2\pi i r x))$. 
			\end{example}
			
			\textbf{Fact} If $f$ is closed injective immersion then $Im(f)$ is an embedded submanifold.
			
			
			\subsection{Sard's Theorem}
			\begin{theorem}
				(Sard) Suppose:
			\end{theorem}
			
			\subsection{Reuglar Level Set}
			\begin{Def} Let $f:M\rightarrow N$ be a smooth map:
				\begin{itemize}
					\item A point $c\in M$ is called a \textbf{regular point} if $df_c:T_p M\rightarrow T_p N$ is surjective;
					
					Elsewise, $c$ is called \textbf{critical point}.
					\item A point $d\in N$ is called a \textbf{regular value} if $\Phi^{-1}(c)$ is a collection of regular points;
					
					Such $\Phi^{-1}(c)$ is called a \textbf{regular value set};
				\end{itemize}
			\end{Def}
			
			\begin{theorem}
				(Regular Value Set Theorem)
			\end{theorem}
			\begin{proof}
				内容...
			\end{proof}
			
		\section{Whitney's Embedding Theorem}
		\section{Whitney's Approxiamation Theorem}
			\begin{Def}
				内容...
			\end{Def}
		\section{Transversality}
		\begin{Def}
			\begin{itemize}
				\item 
				Let: $M$ be a manifold, and $S,S'\subseteq M$ are embedded submanifolds. 	
				
				$S,S'$ are said to \textbf{intersect transversely} if
				\begin{gather*}
					T_p M=T_p S\oplus T_p S'
				\end{gather*}
				
				\item For a smooth map $F:N\rightarrow M$, and $S\subseteq M$ be a submanifold. Then, $F$ is tranverse to $F$ if $T_{F(x)} M=dF_{x }(T_x N)\oplus T_{F(x)} S$. 
			\end{itemize}
		\end{Def}
		
		\begin{proposition}
			内容...
		\end{proposition}
		\begin{proof}
			内容...
		\end{proof}
		
		\begin{corollary}
			内容...
		\end{corollary}
		\begin{proof}
			内容...
		\end{proof}
		\subsection{Graphs}
			\begin{proposition}
				(Global Characterization of Graphs)
			\end{proposition}
			\begin{proof}
				内容...
			\end{proof}
			
			\begin{proposition}
				(Local Characterization of Graphs)
			\end{proposition}
			\begin{proof}
				内容...
			\end{proof}
			
			\begin{proposition}
				(Parametric Transversality Theorem)
			\end{proposition}
			\begin{proof}
				内容...
			\end{proof}
		\subsection{Transversility Homotopy Theorem}
		
		\section{Properties of Morphisms}
	\chapter{Flows}
	\chapter{Differential Forms}
		\section{Differential Form}
			\subsection{Linear Algebras}
			Here is some of review of linear algebra.
			\begin{Def}
				Take the following notions: $k$-multilinear map $T^k(V)$, and symmetric/Alternating $\Lambda^k(V)$.
				
				And
				\begin{gather*}
					Alt(T):=\frac{1}{k!}\sum_{\sigma\in S_k} sgn(\sigma)T^\sigma;T\in T^k(V)\\
					\wedge:\Lambda^k(V)\times \Lambda^l(V)\rightarrow \Lambda^{k+l}(V)\\
					\Lambda(V):=\oplus_{k=0}^n \Lambda_k(V)
				\end{gather*}
			\end{Def}
			\begin{proposition}
				We have
				\begin{gather*}
					\begin{cases}
						dim\Lambda(V)=2^d\\
						dim\Lambda_k(V)=\binom{d}{k},0\leq k\leq d
					\end{cases},dim
					\wedge_k(V)=\{0\};k\geq d
				\end{gather*}
			\end{proposition}
			\begin{proof}
				内容...
			\end{proof}
					\subsubsection{Univeresal Mapping Property}
						\begin{proposition}
							Denote $\varphi:V\times \dots\times V\rightarrow \Lambda_k(V),(v_1,\dots,v_k)\mapsto v_1\wedge\dots\wedge v_k$. Then, for arbitrary alternating multilinear map $h$ of $V\times \dots\times V$ into the vector space $W$. There exists a unique map $\tilde{h}:\wedge_k(V)\rightarrow W$ such that
							\begin{gather*}
								\tilde{h}\circ \varphi=h
							\end{gather*}
							With the following diagram commutes
							\begin{center}
								\begin{tikzcd}
									\Lambda_k(V)\arrow[rd,dashrightarrow]\\
									V\times \dots\times V\arrow[r,"h"]\arrow[u,"\varphi"]&W
								\end{tikzcd}
							\end{center}
						\end{proposition}
						\begin{proof}
							内容...
						\end{proof}
						
						\begin{corollary}
							Suppose that $W=F$. Then, we could take the notion of
							\begin{gather*}
								\wedge_k(V)^*\cong A_k(V)
							\end{gather*}
						\end{corollary}
						\begin{proof}
							内容...
						\end{proof}
					\subsubsection{Pairing}
						\begin{Def}
							Let $V$ and $W$ be real finite dimensional vector spaces, a \textbf{pairing} of $V$ and $W$ is a bilinear map
							\begin{gather*}
								(-,-):V\times W\rightarrow \RR
							\end{gather*}
							A \textbf{pairing} is called \textbf{non-sigular} if whenever $w\not=0$ in $W$, theere exists 
						\end{Def}
					\subsubsection{Linear Transformation of $\Lambda(V)$}
						Take the notion
						\begin{gather*}
							End(\Lambda(V))
						\end{gather*}
						for the endomorphisms of $\Lambda(V)$
			\subsection{Tensor Fields and Differential Forms}
				\subsubsection{Differential Forms}
					\begin{Def}
						Define:
						\begin{gather*}
							\Lambda^k M:=\sqcup_{p\in M} \Lambda^k(T_p M);
						\end{gather*}
						
						The $k$-forms locally written as:
						\begin{gather*}
							\omega=\sum_I \omega^I dx^I\qquad dx^I=dx^{i_1}\wedge\dots\wedge dx^{i_k}\qquad 1\leq i_1<\dots<i_k\leq n
						\end{gather*}
						
						Take the following notion $\Omega^k(M)$
					\end{Def}
					
					And action
					\begin{gather*}
						dx^{i_1}\wedge dx^{i_k}(\frac{\partial}{\partial x_{j_1}},\dots,\frac{\partial}{\partial x^{i_k}})=\delta^I_J
					\end{gather*}
					
			\subsection{Exterior Derivative}
				\begin{Def}
					Define: For $k$-form $\omega=\sum_J \omega_J dx^J dx^J$, then 
					\begin{gather*}
						d\omega:=\sum d\omega_J\wedge dx^J\\
						\qquad=\sum_J \sum_{i} \frac{\partial_J}{\partial x^i}dx^i\wedge dx^J
					\end{gather*}
				\end{Def}
				
				\begin{proposition}
					\begin{itemize}
						\item $d$ is linear over $\RR$;
						\item Given $k$-form $\omega$, $l$-form $\eta$, $U\subseteq \RR^n$ or $\mathbb{H}^n$
						\item $d^2=0$;
						\item $d$ commutes with pullback
						\begin{gather*}
							F^*(d\omega)=d(F^*\omega)
						\end{gather*}
					\end{itemize}
				\end{proposition}
				\begin{proof}
					内容...
				\end{proof}
				
				\begin{theorem}
					(Existence and Uniqueness of Exterior Differentation)
				\end{theorem}
				\begin{proof}
					内容...
				\end{proof}
				
				\subsubsection{Closed,Exact Form}
				\begin{Def}
					Consider the cohomology: $\dots\rightarrow H^{k+1}_{dR}(M)\rightarrow H^k_{dR}(M)\rightarrow\dots$. Then, the differential form $\omega\in \Omega(M)$ is:
					\begin{itemize}
						\item Closed, if $d\omega=0$;
						\item Exact, if there exsits some $f$ such that $\omega=df$;
					\end{itemize}
				\end{Def}
				
				Especially, for $1$-form:
				\begin{proposition}
					For $1$-form $\omega=\sum a_i dx^i$, we have
					\begin{gather*}
						d\omega=\sum (\frac{\partial a_i}{\partial x_i}-\frac{\partial a_j}{x_j})dx^i\wedge dx^j
					\end{gather*}
					
					Thus closedness$\Leftrightarrow \frac{\partial a_i}{\partial x_i}=\frac{\partial a_j}{\partial x_j}$
				\end{proposition}
			\subsection{The Lie Derivative}
		\section{Orientation}
			\subsection{Orientation of Vector Spaces}
				\begin{Def}
					Let $V$ be an $N$-dimensional real inner product space. We extend the inner product from $V$ to all of $\Lambda(V)$
					\begin{gather*}
						<w_1\wedge \dots\wedge w_p,v_1\wedge\dots\wedge v_p>:=det<w_i,v_j>
					\end{gather*}
					The bilinear map can be extended linearly to $\Lambda_p(V)$.
				\end{Def}
				
				\begin{proposition}
					If $e_1,\dots,e_n$ is an orthonormal basis of $V$, then the corresponding basis of $\Lambda(V)$ is an orthonormal basis for $\Lambda(V)$.
				\end{proposition}
				\begin{proof}
					内容...
				\end{proof}
				
				
				\begin{Def}
					(Consistently Oriented) Assume that $(v_1,\dots,v_n)$ and $(w_1,\dots,w_n)$ be two bases of $V$ with $dim(V)=n$, then we have the transform matrix
					\begin{gather*}
						v^T=Bw^T
					\end{gather*}
					Then, they are \textbf{consistently oriented} if $det(B)>0$.
				\end{Def}
			\subsection{Orientation of Manifold}
				\begin{Def}
					Let $M$ be a manifold. Then a frame
					\begin{gather*}
						(E_1,\dots,E_n)
					\end{gather*}
					is \textbf{(positively) oriented} if
				\end{Def}
			\subsection{The Orientation Covering}
			\subsection{The Orientation of Hypersurface}
			\subsection{The Orientation of }
	\chapter{Integration on Manifolds}
		\section{Integration on Manifolds}
		\subsection{Euclidean Space}
			\begin{Def}
				Defien:
				\begin{gather*}
					f:D\rightarrow \RR\qquad \int_D fdx^1\wedge\dots\wedge dx^n:=\int_D fdx^1\dots dx^n
				\end{gather*}
			\end{Def}
		\subsection{Manifolds}
		\begin{Def}
			Let $M$ be a manifold. A \textbf{partition of unity} is a collection of smooth maps $(\varphi_i:M\rightarrow \RR)$ 
			\begin{itemize}
				\item there exists $\varphi_i:U_i\rightarrow V_i$ such that
				\begin{gather*}
					supp(\varphi_i)\subseteq U_i
				\end{gather*}
				
				\item For each $p\in M$, only finitely many of $\varphi_i(p)\not=0$.
			\end{itemize}
		\end{Def}
		\begin{Def}
			Let $M$ be an oriented n-dim smooth manifold, $\omega$ a diff n-form, local integrable. Choose a partition of unity $(\varphi_i)$, define
			\begin{gather*}
				\int \omega=\sum_i \int h_i^{-1*}(\varphi_i\circ\omega) 
			\end{gather*}
		\end{Def}
	\chapter{De Rham Cohomology}
		\section{De Rham Complex}
	\chapter{Integratability}
		\section{Integral Curve}
			\begin{Def}
				Let: $M$ be a smooth manifold, $X\in \mathfrak{X}(M)$. 
				
				A \textbf{integral curve} is a smooth curve $\gamma:I\rightarrow M$ such that, for every $t\in I$,
				\begin{gather*}
					\dot{\gamma}(t)=X_{\gamma(t)}
				\end{gather*}
				
				where $\dot{\gamma}(t)=(d\gamma)_t(\frac{\partial}{\partial t})\in T_{\gamma(t)}M$
			\end{Def}
			
			Sometimnes we will use the notion $d_{\gamma_s}(1)$, because $1\in \RR$ corresponds the standard tangent vector $\frac{\partial}{\partial t}|s$.
			
			Or, another explaination is to consider the curve goes through $t$: $c(r)=t+r$, $d\gamma_s(1)=d\gamma_s(c'(0))=(\gamma\circ c)'(0)=\frac{d}{dr}|_{r=0} \gamma(s+r)$.
		
			\subsection{Existence of Maximal Integral Curve}
					\begin{proposition}
					(Frame Coordinate of Integral Curve) Take frame bundle $(x^1,\dots,x^n$), and $X\in \mathfrak{X}(M)$, then the equation of integral curve at $p$ with coordinate $\varphi=(x^1,\dots,x^n):U\rightarrow \RR^n$, then
					\begin{gather*}
						X|U=\sum X^i(x)\frac{\partial}{\partial x^i}\\
						x(t)=\varphi(\gamma(t))
					\end{gather*}
					
					We obtain
					\begin{gather*}
						\dot{x(t)}=(X^1(x(t)),\dots,X^n(x(t)))
					\end{gather*}
					
					Especially, it is local-Lipschitz. 
				\end{proposition}
				\begin{proof}
					内容...
				\end{proof}
				\begin{proposition}
					(Local Existence) For each point $p\in M$, there exists a local integral curve $\gamma$ which goes through $p$.
				\end{proposition}
				\begin{proof}
					Apply Picard-Lindelof.
				\end{proof}
			\subsection{Maximal Integral Curve}
				\begin{Def}
					Let $X\in \mathfrak{X}(M)$. Now, consider the inclusion of curves:
					\begin{gather*}
						(\gamma:I\rightarrow M,I)\subseteq (\theta:J\rightarrow M,J)\Leftrightarrow i\subseteq J,\theta|I=\gamma
					\end{gather*}
				\end{Def}
				
				\begin{proposition}
					There exists a unique maximal curve which goes through $p$. (We could take $\gamma_p(0)=p$)
				\end{proposition}
				\begin{proof}
					By uniqueness of ODE. 
				\end{proof}
			\subsection{Visualization of Integral Curve}
			\subsection{Compelete Vector Field}
				\begin{Def}
					A vector field $X\in \mathfrak{X}(M)$ is called \textbf{compelete}, if for each $p\in M$, there exists a \underline{unique} maximal integral curve
					\begin{gather*}
						\gamma_p:I_p\rightarrow M,\gamma_p(0)=p\qquad I_p=\RR,\forall p\in M
					\end{gather*}
				\end{Def}
		\section{Distributions and Involutivity}
		\section{Frobenius Theorems}
		
	\chapter{Appendix: Smooth Category}
		\begin{Def}
			内容...
		\end{Def}
\end{document}